% Written By Enoch Yu
% Downloaded from archive.enochyu.com

\documentclass{article}

\usepackage{amsmath}
\usepackage{amssymb}
\usepackage{amsthm}
\usepackage{mathtools}

\usepackage[margin=1in, headsep=15pt]{geometry}
\usepackage{lastpage}
\usepackage{fancyhdr}
\fancyhead[R]{Enoch Yu}
\fancyfoot[C]{Page \thepage \hspace{1pt} of \pageref{LastPage}}
\pagestyle{fancy}
\usepackage{parskip}

\usepackage{enumitem}
\usepackage[dvipsnames]{xcolor}
\usepackage[normalem]{ulem}
\usepackage{url}
\usepackage{hyperref}
\hypersetup{
    colorlinks=true,
    linkcolor=black,
    filecolor=magenta,      
    urlcolor=cyan,
    pdftitle={7.22.2025 AMC 10,12 Hard Problems},
    pdfpagemode=FullScreen,
}
\usepackage{tabularx}
\usepackage{arydshln}
\usepackage{float}
\usepackage{pdfpages}

\usepackage{tikz}
\usetikzlibrary{calc, angles, shapes.geometric}
\usepackage{tkz-euclide}
\usepackage{pgfplots}
\pgfplotsset{compat=1.9}
\usepackage[font=small,labelfont=bf]{caption}

\newtheorem{theorem}{Theorem}[section]
\newtheorem{lemma}[theorem]{Lemma}
\newtheorem*{lemma*}{Lemma}
\newtheorem{sublemma}{Lemma}[section]
\newtheorem{proposition}{Proposition}
\newtheorem{corollary}{Corollary}[theorem]
\newtheorem{example}{Example}[section]
\newtheorem*{example*}{Example}
\newenvironment{solution}{\begin{trivlist}\item[]{\bf Solution}}{\qed \end{trivlist}}


\title{7.22.2025 AMC 10,12 Hard Problems}
\author{Enoch Yu}
\date{22 July 2025}

\begin{document}
\subsection*{Problem 1.}
{\color {BrickRed}
Let $P(x)$ be a nonzero polynomial such that
$(x - 1)P(x + 1) = (x + 2)P(x)$ for every real $x$,
and $P(2)^2 = P(3)$. Then $P \left( \frac{7}{2} \right)
= \frac{m}{n}$, where $m$ and $n$ are relatively prime
positive integers. Find $m + n$.
}

\begin{solution}
Substitute $x = -1, 0, 1$. Substituting the values
leads to the fact that $P(-1) = P(0) = P(1) = 0$.
Therefore, let
\[
P(x) = f(x) (x + 1) (x) (x - 1).
\]
Substituting to the equation and assuming that
$x \neq -1, 0, 1$ leads to the following equations.
\begin{align*}
    (x - 1) f(x + 1) (x + 2) (x + 1) (x)
    &= (x + 2) f(x) (x + 1) (x) (x - 1) \\
    f(x + 1) &= f(x)
\end{align*}
In other words, $f(x)$ is a constant for $x \neq -1, 0, 1$.
\begin{align*}
    \left( f(2) \cdot 3 \cdot 2 \cdot 1 \right)^2
    &= f(3) \cdot 4 \cdot 3 \cdot 2 \\
    f(2)^2 \cdot 36 &= f(2) \cdot 24 \\
    f(2) &= \frac{2}{3} \quad (\because f(2) \neq 0) 
\end{align*}
Substituting $x = \frac{7}{2}$ to $P(x)$,
$P \left( \frac{7}{2} \right)
= f\left( \frac{7}{2} \right)
    \left( \frac{7}{2} + 1 \right)
    \left( \frac{7}{2} \right)
    \left( \frac{7}{2} - 1 \right)
= f(2) \cdot \frac{315}{8}$, which is $\frac{105}{4}$.
Thus, $105 + 4 = \boxed{109}$.
\end{solution}

\subsection*{Problem 2.}
{\color {BrickRed}
For certain real number $a$, $b$, and $c$, the polynomial
\[
g(x) = x^3 + ax^2 + x + 10
\]
has three distinct roots, and each root of $g(x)$ is also
a root of the polynomial
\[
f(x) = x^4 + x^3 + bx^2 + 100x + c
\]
What is $f(1)$?
}

\begin{solution}
$f(x)$ could be rewritten as $(x - m)(g(x))$.
The coefficients of $x^3$, $x^2$, $x$, and $x^0$
could be compared.
\begin{align*}
    f(x)
    &= (x - m)(x^3 + ax^2 + x + 10) \\
    &= x^4 + (a - m)x^3 + (1 - am)x^2 + (10 - m)x - 10m
\end{align*}
Therefore, $m = -90$, $a = -89$, $b = -8009$, and
$c = 900$. Thus, $f(1) = 1 + 1 + b + 100 + c = \boxed{-7007}$.
\end{solution}

\subsection*{Problem 3.}
{\color {BrickRed}
Let $S$ be the number of ordered pairs of integers
$(a, b)$ with $1 \leq a \leq 100$ and $b \geq 0$ such
that the polynomial $x^2 + ax + b$ can be factored
into the product of two (not necessarily distinct)
linear factors with integer coefficients. Find the
remainder when $S$ is divided by $1000$.
}

\begin{solution}
Let $-p$ and $-q$ be the roots of the polynomial
$x^2 + ax + b$. In other words, $p + q = a$ and $pq = b$.
Since the equations are symmetric, WLOG, let $p \geq q$.
\[
\begin{array}{ccc}
    p & q & \text{Number of Cases} \\
    0 & 1 \sim 100 & 100 \\
    1 & 1 \sim 99 & 99 \\
    2 & 2 \sim 98 & 97 \\
    3 & 3 \sim 97 & 95 \\
    \vdots & \vdots & \vdots \\
    50 & 50 & 1
\end{array}
\]
Because an equation with two roots $p$ and $q$ is unique,
all cases are counted. Thus, the total number of cases
are $1 + 3 + \cdots + 97 + 99 + 100
= \frac{100 \cdot 50}{2} + 100 = 2600$. Therefore,
the answer is $\boxed{600}$.

\end{solution}

\subsection*{Problem 4.}
{\color {BrickRed}
The polynomial $P(x)$ is cubic. What is the largest value
of $k$ for which the polynomials
\[
Q_1(x) = x^2 + (k - 29)x - k \quad \text{and}
    \quad Q_2(x) = 2x^2 + (2k - 43)x + k
\]
are both factors of $P(x)$?
}

\begin{solution}
Notice that because $k$ is a constant, $P(x)$ could be rewritten.
\begin{align*}
    P(x)
    &= Q_1(2ax + b) = (x^2 + (k - 29)x - k)(2ax + b) \\
    &= Q_2(ax - b) = (2x^2 + (2k - 43)x + k)(ax - b)
\end{align*}
Therefore, the coefficients of $x^2$ and $x$ could be
compared for the value of $k$.
\begin{align*}
    2ak - 58a + b &= 2ak - 43a - 2b \\
    b &= 5a
    \\\\
    kb - 29b - 2ak &= -2bk + 43b + ak \\
    3bk &= 3ak + 72b \\
    15ak &= 3ak + 360a \\
    \therefore k &= \boxed{30}
\end{align*}
\end{solution}

\subsection*{Problem 5.}
{\color {BrickRed}
A real number $a$ is chosen randomly and uniformly
from the interval $[-20, 18]$. Calculate the probability
that the roots of the polynomial
\[
x^4 + 2a x^3 + (2a - 2)x^2 + (-4a + 3)x - 2
\]
are all real.
}

\begin{solution}
Notice that when $x = 1$ and $x = -2$, the polynomial
is zero. Thus, with synthetic division, the polynomial
could be factored.
\[
(x - 1)(x + 2)(x^2 + (2a - 1)x + 1) = 0
\]
For all the roots of $x$ to be real, $(2a-1)^2 - 4$
must be greater than or equal to zero.
\begin{align*}
    4a^2 - 4a - 3 &\geq 0 \\
    (2a + 1)(2a - 3) &\geq 0 \\
    -20 \leq a \leq -\frac{1}{2}, 18 \geq a \geq \frac{3}{2}
\end{align*}
Lengths could be used to compute the probability.
\begin{align*}
    \frac{\left( 18 - \frac{3}{2} \right)
        + \left( 20 - \frac{1}{2} \right)}{18 + 20}
    &= \frac{\frac{33}{2} + \frac{39}{2}}{38} \\
    &= \frac{36}{38} \\
    &= \boxed{\frac{18}{19}}
\end{align*}
\end{solution}

\subsection*{Problem 6.}
{\color {BrickRed}
For what real values of $k$ do
\[
1988x^2 + kx + 8891 \quad \text{and} \quad 8891x^2 + kx + 1988
\]
have a common zero?
}

\begin{solution}
Let $\alpha$ be the common root.
\begin{align*}
    1988\alpha^2 + k\alpha + 8891 &= 0 \\
    8891\alpha^2 + k\alpha + 1988 &= 0
\end{align*}
Subtracting two equations, $6903\alpha^2 = 6903$ is obtained.
In other words, the common root could be either $1$ or $-1$.
Therefore, $k$ could be $\boxed{\pm 10879}$.
\end{solution}
\end{document}

